\documentclass[../main.tex]{subfiles}

\begin{document}

\ifSubfilesClassLoaded{

    %\tableofcontents
    
    %\setcounter{chapter}{}
}{}

\chapter{ Submersions, Immersions and Embedings }
\subsection{Local Diffeomorphism}
\begin{definition}{}{}
    If $M$ and $N$ are smooth manifolds with or without boundary, a map $F: M \to N$ is called a \dfntxt{local diffeomorphism} if every point $p \in M$ has a neighborhood $U$ such that $F(U)$ is open in $N$ and $F|_U : U \to F(U)$ is a diffeomorphism. 
\end{definition}
\begin{theorem}{Inverse Function Theorem for Manifolds}{}
Suppose $M$ and $N$ are smooth manifolds, and $F: M \to N$ is a smooth map. If $p \in M$ is a point such that $dF_p$ is invertible, then there are connected neighborhoods $U_0$ of $p$ and $V_0$ of $F(p)$ such that $F|_{U_0}: U_0 \to V_0$ is a diffeomorphism.
\end{theorem}

\begin{note}
    也就是说, 光滑映射在微分映射可逆的点附近是微分同胚.
\end{note}

\end{document}